%package list
\documentclass{article}
\usepackage[top=3cm, bottom=3cm, outer=3cm, inner=3cm]{geometry}
\usepackage{multicol}
\usepackage{graphicx}
\usepackage{url}
%\usepackage{cite}
\usepackage{hyperref}
\usepackage{array}
%\usepackage{multicol}
\newcolumntype{x}[1]{>{\centering\arraybackslash\hspace{0pt}}p{#1}}
\usepackage{natbib}
\usepackage{pdfpages}
\usepackage{multirow}
\usepackage[normalem]{ulem}
\useunder{\uline}{\ul}{}
\usepackage{svg}
\usepackage{xcolor}
\usepackage{listings}
\lstdefinestyle{ascii-tree}{
    literate={├}{|}1 {─}{--}1 {└}{+}1 
  }
\lstset{basicstyle=\ttfamily,
  showstringspaces=false,
  commentstyle=\color{red},
  keywordstyle=\color{blue}
}
%\usepackage{booktabs}
\usepackage{caption}
\usepackage{subcaption}
\usepackage{float}
\usepackage{array}

\newcolumntype{M}[1]{>{\centering\arraybackslash}m{#1}}
\newcolumntype{N}{@{}m{0pt}@{}}


%%%%%%%%%%%%%%%%%%%%%%%%%%%%%%%%%%%%%%%%%%%%%%%%%%%%%%%%%%%%%%%%%%%%%%%%%%%%
%%%%%%%%%%%%%%%%%%%%%%%%%%%%%%%%%%%%%%%%%%%%%%%%%%%%%%%%%%%%%%%%%%%%%%%%%%%%
\newcommand{\itemEmail}{vmamanian@unsa.edu.pe}
\newcommand{\itemStudent}{Victor Mamani Anahua}
\newcommand{\itemCourse}{Fundamentos de la Programación II}
\newcommand{\itemCourseCode}{20230489}
\newcommand{\itemSemester}{II}
\newcommand{\itemUniversity}{Universidad Nacional de San Agustín de Arequipa}
\newcommand{\itemFaculty}{Facultad de Ingeniería de Producción y Servicios}
\newcommand{\itemDepartment}{Departamento Académico de Ingeniería de Sistemas e Informática}
\newcommand{\itemSchool}{Escuela Profesional de Ingeniería de Sistemas}
\newcommand{\itemAcademic}{2023 - B}
\newcommand{\itemInput}{Del 15 Enero 2024}
\newcommand{\itemOutput}{Al 22 Enero 2024}
\newcommand{\itemPracticeNumber}{22}
\newcommand{\itemTheme}{Laboratorio 22}
%%%%%%%%%%%%%%%%%%%%%%%%%%%%%%%%%%%%%%%%%%%%%%%%%%%%%%%%%%%%%%%%%%%%%%%%%%%%
%%%%%%%%%%%%%%%%%%%%%%%%%%%%%%%%%%%%%%%%%%%%%%%%%%%%%%%%%%%%%%%%%%%%%%%%%%%%

\usepackage[english,spanish]{babel}
\usepackage[utf8]{inputenc}
\AtBeginDocument{\selectlanguage{spanish}}
\renewcommand{\figurename}{Figura}
\renewcommand{\refname}{Referencias}
\renewcommand{\tablename}{Tabla} %esto no funciona cuando se usa babel
\AtBeginDocument{%
	\renewcommand\tablename{Tabla}
}

\usepackage{fancyhdr}
\pagestyle{fancy}
\fancyhf{}
\setlength{\headheight}{30pt}
\renewcommand{\headrulewidth}{1pt}
\renewcommand{\footrulewidth}{1pt}
\fancyhead[L]{\raisebox{-0.2\height}{\includegraphics[width=3cm]{img/logo_episunsa.png}}}
\fancyhead[C]{\fontsize{7}{7}\selectfont	\itemUniversity \\ \itemFaculty \\ \itemDepartment \\ \itemSchool \\ \textbf{\itemCourse}}
\fancyhead[R]{\raisebox{-0.2\height}{\includegraphics[width=1.2cm]{img/logo_abet}}}
\fancyfoot[L]{Estudiante Victor Mamani A.}
\fancyfoot[C]{\itemCourse}
\fancyfoot[R]{Página \thepage}

% para el codigo fuente
\usepackage{listings}
\usepackage{color, colortbl}
\definecolor{dkgreen}{rgb}{0,0.6,0}
\definecolor{gray}{rgb}{0.5,0.5,0.5}
\definecolor{mauve}{rgb}{0.58,0,0.82}
\definecolor{codebackground}{rgb}{0.95, 0.95, 0.92}
\definecolor{tablebackground}{rgb}{0.8, 0, 0}

\lstdefinestyle{java}{frame=tb,
	language=Java,
	showstringspaces=false,
	columns=flexible,
	basicstyle={\footnotesize\ttfamily\color[RGB]{255,255,255}},
	numberstyle=\color{mygray},
	numbers=left, 
	keywordstyle=\color{myblue},
	morekeywords={String, System},
	commentstyle=\color{mygray},
	stringstyle=\color{mygreen},
	breaklines=true,
	breakatwhitespace=true,
	tabsize=2,
	backgroundcolor= \color{codebackgroundCode},
	showspaces=false,
	showtabs=false,
	showlines=false,
}

\lstset{frame=tb,
	language=bash,
	aboveskip=3mm,
	belowskip=3mm,
	showstringspaces=false,
	columns=flexible,
	basicstyle={\small\ttfamily},
	numbers=none,
	numberstyle=\tiny\color{gray},
	keywordstyle=\color{blue},
	commentstyle=\color{dkgreen},
	stringstyle=\color{mauve},
	breaklines=true,
	breakatwhitespace=true,
	tabsize=3,
	backgroundcolor= \color{codebackground},
}

\begin{document}
	
	\vspace*{10px}
	
	\begin{center}	
		\fontsize{17}{17} \textbf{ Informe de Laboratorio \itemPracticeNumber}
	\end{center}
	\centerline{\textbf{\Large Tema: \itemTheme}}
	%\vspace*{0.5cm}	

	\begin{flushright}
		\begin{tabular}{|M{2.5cm}|N|}
			\hline 
			\rowcolor{tablebackground}
			\color{white} \textbf{Nota}  \\
			\hline 
			     \\[30pt]
			\hline 			
		\end{tabular}
	\end{flushright}	

	\begin{table}[H]
		\begin{tabular}{|x{4.7cm}|x{4.8cm}|x{4.8cm}|}
			\hline 
			\rowcolor{tablebackground}
			\color{white} \textbf{Estudiante} & \color{white}\textbf{Escuela}  & \color{white}\textbf{Asignatura}   \\
			\hline 
			{\itemStudent \par \itemEmail} & \itemSchool & {\itemCourse \par Semestre: \itemSemester \par Código: \itemCourseCode}     \\
			\hline 			
		\end{tabular}
	\end{table}		
	
	\begin{table}[H]
		\begin{tabular}{|x{4.7cm}|x{4.8cm}|x{4.8cm}|}
			\hline 
			\rowcolor{tablebackground}
			\color{white}\textbf{Laboratorio} & \color{white}\textbf{Tema}  & \color{white}\textbf{Duración}   \\
			\hline 
			\itemPracticeNumber & \itemTheme & 04 horas   \\
			\hline 
		\end{tabular}
	\end{table}
	
	\begin{table}[H]
		\begin{tabular}{|x{4.7cm}|x{4.8cm}|x{4.8cm}|}
			\hline 
			\rowcolor{tablebackground}
			\color{white}\textbf{Semestre académico} & \color{white}\textbf{Fecha de inicio}  & \color{white}\textbf{Fecha de entrega}   \\
			\hline 
			\itemAcademic & \itemInput &  \itemOutput  \\
			\hline 
		\end{tabular}
	\end{table}
	
	\section{Tarea}
	\begin{itemize}		
		\item Cree una versión del videojuego de estrategia usando componentes básicos GUI: Etiquetas, botones,
		cuadros de texto, JOptionPane, Color.
		\item Además, utilizar componentes avanzados GUI: Layouts, JPanel, áreas de texto, checkbox, botones de
		radio y combobox.
		\item Considerar nivel estratégico y táctico.
		\item Considerar hasta las unidades especiales de los reinos.
		\item Hacerlo iterativo.
	\end{itemize}

	\section{Equipos, materiales y temas utilizados}
	\begin{itemize}
		\item Sistema Operativo Ubuntu GNU Linux 23 lunar 64 bits Kernell 6.2.v
		\item Visual Studio Code.
		\item VIM 9.0.
		\item OpenJDK 64-Bits 19.0.7.
		\item Git 2.39.2.
		\item Cuenta en GitHub con el correo institucional.
		\item Programación Orientada a Objetos.
		\item Actividades del Laboratorio 22.	
	\end{itemize}
	
	\section{URL de Repositorio Github}
	\begin{itemize}
		\item URL del Repositorio GitHub para clonar o recuperar.
		\item \url{https://github.com/VictorMA18/fp2-23b.git}
		\item URL para el laboratorio 22 en el Repositorio GitHub.
		\item \url{https://github.com/VictorMA18/fp2-23b/tree/main/Fase03/Lab22}
	\end{itemize}
	
	\section{Actividades del Laboratorio 22}

	\subsection{Ejercicio Soldado(Herencias)}
	\begin{itemize}	
		\item En esta seccion solo reutilizamos las Herencias de la clase Soldado.
		\item El codigo y el commit seria el siguiente:
	\end{itemize}	
	\begin{lstlisting}[language=bash,caption={Commit}][H]
		$ git commit -m "Agregando la clase Soldado y Espadachin para poder hacer el juego bueno solo en la clase espadachin usamos la herencia que nos deja la clase Soldado y tambien creamos la funcion muroEscudo() la cual devuelve como mensaje el uso de esta habilidad defensiva y los getters y setters"
	\end{lstlisting}	
	\begin{lstlisting}[language=java,caption={Las lineas de codigos de la clase Espadachin creada:}][H]
		public class Espadachin extends Soldado{
			private int swordlth;
			public Espadachin(String name , int attacklevel, int defenselevel, int lifelevel, int speed, String attitude ,boolean lives, int row, String column, int swordlth){
				super(name, attacklevel, defenselevel, lifelevel, speed, attitude, lives, row, column);
				this.swordlth = swordlth;
			}
			public void muroEscudo(){
				System.out.println("Usted uso la habilidad muro de Escudos");
			}
			public int getSwordlth(){
				return swordlth;
			}
			public void setSwordlth(int n){
				this.swordlth = n;
			}
		}
	\end{lstlisting}
	\begin{lstlisting}[language=java,caption={Las lineas de codigos de la clase Caballero creada:}][H]
		public class Caballero extends Soldado{
			private boolean montar;
			private String arma;
			public Caballero(){
			}
			public Caballero(String name , int attacklevel, int defenselevel, int lifelevel, int speed, String attitude ,boolean lives, int row, String column,boolean montar){
				super(name, attacklevel, defenselevel, lifelevel, speed, attitude, lives, row, column);
				this.montar = montar;
			}
			public void montar(){
				if(!this.montar){
					this.arma = "Lanza";
					this.embestir();
				}
			}
			public void desmontar(){
				if (this.montar) {
					this.arma = "Espada";
				}
			}
			public void embestir(){
				if(!montar){
					this.atacar();
					this.atacar();
				}else{
					this.atacar();
					this.atacar();
					this.atacar();
				}
			}
			public String getArma(){
				return arma;
			}
		}
	\end{lstlisting}
	\begin{lstlisting}[language=java,caption={Las lineas de codigos de la clase Arquero creada:}][H]
		public class Arquero extends Soldado{
			private int flechas;
			public Arquero(){
				
			}
			public Arquero(String name , int attacklevel, int defenselevel, int lifelevel, int speed, String attitude ,boolean lives, int row, String column, int flechas){
				super(name, attacklevel, defenselevel, lifelevel, speed, attitude, lives, row, column);
				this.flechas = flechas;
			}
			public void disparar(){
				if(this.flechas == 0){
					System.out.println("El arquero ya tiene flechas para poder disparar");
				}else{
					this.flechas = flechas - 1;
					this.atacar();
				}
			}
			public void setFlechas(int n){
				this.flechas = n;
			}
			public int getFlechas(){
				return flechas;
			}
		}
	\end{lstlisting}
	\begin{lstlisting}[language=java,caption={Las lineas de codigos de la clase Lancero creada:}][H]
		public class Lancero extends Soldado{
			private int lancelth;
			public Lancero(){
			}
			public Lancero(String name , int attacklevel, int defenselevel, int lifelevel, int speed, String attitude ,boolean lives, int row, String column, int lancelth){
				super(name, attacklevel, defenselevel, lifelevel, speed, attitude, lives, row, column);
				this.lancelth = lancelth;
			}
			public void schiltrom(){
				this.setDefenseLevel(this.getDefenseLevel() + 1);
				System.out.println("El lancero uso el schiltrom su nivel de defensa subio 1 punto");
			}
			public void setLancelth(int n){
				this.lancelth = n;
			}
			public int getLancelth(){
				return lancelth;
			}
		}
	\end{lstlisting}
	\begin{lstlisting}[language=java,caption={Las lineas de codigos de la clase Soldado creada:}][H]
		// Laboratorio Nro 22 - Ejercicio Soldado
		// Autor: Mamani Anahua Victor Narciso
		// Colaboro:
		// Tiempo:
		import java.util.*;
		public class Soldado { //CREAMOS LA CLASE SOLDODADO PARA PODER USAR UN ARREGLO BIDIMENSIONAL DONDE NECESITAMOS LA VIDA , EL NOMBRE DEL SOLDADO Y TAMBIEN SU POSICION COMO LA FILA Y LA COLUMNA   
		
			private String name;
			private int lifeactual;
			private int row;
			private String column;
			private int attacklevel;
			private int defenselevel;
			private int lifelevel;
			private int speed;
			private String attitude;
			private boolean lives;
		
			Random rdm = new Random();
		
			//Anadiendo metodo que nos permita que un arreglo tenga datos nulos si este esta vacio
			public Soldado(){
				this.name = "";
				this.row = 0;
				this.column  = "";
				this.attacklevel = 0;
				this.defenselevel = 0;
				this.lifelevel = 0;
				this.lifeactual = 0;
				this.speed = 0;
				this.attitude = "";
				this.lives = false;
			}
		
			//Constructor
			public Soldado(String name, int health, int row, String column){
				this.name = name;
				this.lifeactual = health;
				this.lifelevel = health;
				this.lifeactual = health;
				this.row = row;
				this.column = column;
				this.lives = true;
				
				//YA QUE ESTOS DATOS SERIAN ALEATORIOS YA QUE SE ESTARIA CREANDO EL SOLDADO TENDRIAMOS DATOS QUE SERIAN COMO ATTACKLEVEL DEFENSELEVEL EL CUAL TENDRIAN QUE SER ALEATORIOS    
				this.attacklevel = rdm.nextInt(5) + 1;
				this.defenselevel = rdm.nextInt(5) + 1;
		
			}
			
			//Constructor para los diferentes niveles como de vidad defensa ataque velocidad
			public Soldado(String name , int attacklevel, int defenselevel, int lifelevel, int speed, String attitude ,boolean lives, int row, String column) {
				this.name = name;
				this.attacklevel = attacklevel;
				this.defenselevel = defenselevel;
				this.lifeactual = lifelevel;
				this.lifelevel = lifelevel;
				this.speed = speed;
				this.lives = lives;
				this.row = row;
				this.column = column;
				this.attitude = attitude;
			}
		
			//Metodos necesarios como avanzar defender huir al seratacado al retroceder
			public void advance(){
				this.speed = getSpeed() + 1;
				System.out.println("El soldado " + this.name + "avanzo");
			}
			public void defense(){
				this.speed = 0;
				this.attitude = "DEFENSIVA";
				System.out.println("El soldado " + this.name + "esta defendiendo");
			}
			public void flee(){
				this.speed = getSpeed() + 2;
				this.attitude = "HUYE";
				System.out.println("El soldado " + this.name + "esta huyendo");
			}
			public void back(){
				System.out.println("El soldado " + this.name + "esta retrocediendo");
				if(this.speed == 0){
					this.speed = rdm.nextInt(5) - 5;
				}else{
					if(this.speed > 0){
						this.speed = 0;
						this.attitude = "DEFENSIVA";
					}
				}
			}
			public void atacar(){
				this.speed += 1;
				this.attitude = "Atacar";
				this.lifeactual += 1;
				if(this.lifeactual == 0){
					morir();
				}
			}
			public void attack(Soldado soldier){
				if(this.getLifeActual() > soldier.getLifeActual()){
					int life = this.getLifeActual() - soldier.getLifeActual();
					this.setLifeActual(life);
					this.setLifeLevel(life);
					soldier.lives = false;
					soldier.morir();
					System.out.println(this.name + " asesino al soldado " + soldier.name);
				}else if(soldier.getLifeActual() > this.getLifeActual()){
					int life = soldier.getLifeActual() - this.getLifeActual();
					this.lives = false;
					this.morir();
					soldier.setLifeActual(life);
					soldier.setLifeLevel(life);
					System.out.println(soldier.name + " asesino al soldado " + this.name);
				}else{
					this.lives = false;
					this.morir();   
					soldier.lives = false;
					soldier.morir();
					System.out.println("los 2 soldados se asesinaron");
				}
			}
			public void morir(){
				this.lives = false;
				this.attitude = "SOLDADO MUERTO";
			}
		
			// Metodos mutadores
			public void setName(String n){
				name = n;
			}
			public void setLifeActual(int p){
				lifeactual = p;
			}
			public void setRow(int b){
				row = b;
			}
			public void setColumn(String c){
				column = c; 
			}
			public void setAttackLevel(int attacklevel) {
				this.attacklevel = attacklevel;
			}
			public void setDefenseLevel(int defenselevel) {
				this.defenselevel = defenselevel;
			}
			public void setLifeLevel(int lifelevel){
				this.lifelevel = lifelevel;
			}
			public void setSpeed(int speed) {
				this.speed = speed;
			}
			public void setAttitude(String attitude) {
				this.attitude = attitude;
			}
			public void setLives(boolean lives) {
				this.lives = lives;
			}
		
			// Metodos accesores
			public String getName(){
				return name;
			}
			public int getLifeActual(){
				return lifeactual;
			}
			public int getRow(){
				return row;
			}
			public String getColumn(){
				return column;
			}
			public int getAttackLevel() {
				return attacklevel;
			}
			public int getDefenseLevel() {
				return defenselevel;
			}
			public int getLifeLevel(){
				return lifelevel;
			}
			public int getSpeed() {
				return speed;
			}
			public String getAttitude() {
				return attitude;
			}
			public boolean getLives() {
				return lives;
			}
		
			// Completar con otros metodos necesarios
			public String toString(){ //CREAMOS ESTE METODO PARA IMPRIMIR LOS DATOS DEl OBJETO
				String join = "\nNombre: " + getName() + "\nVida: " + getLifeActual() + "\nFila: " + getRow() + "\nColumna: " + getColumn() + "\nNivel de ataque: " + getAttackLevel() + "\nNivel de Defensa: " + getDefenseLevel() + "\nNivel de vida: " + getLifeLevel() + "\nVelocidad: " + getSpeed() + "\nActitud: " + getAttitude() + "\nEstado: " + getLives(); //Agregamos un espaciador para poder separar
				return join;
			}
			
		}
	\end{lstlisting}

	\subsection{Ejercicio Ejercito y Mapa}
	\begin{itemize}	
		\item En esta seccion En la clase Ejercito creamos un contructor null y constructor el cual reciba parametros para reconocer a este le agregamos los metodos viewsoldiers() , longerlife() , averagelife(),  rankingInsercionLife() y tambien sus getters y setters esto con el fin que nos ayude a simplificar algunas funciones las cuales sean representadas mejor con una clase Ejercito.
		\item Por parte de clase Mapa integramos 2 objetos de la clase Ejercito las cuales nos ayudan con sus metodos para ver la informacion de la batalla entre los 2 ejercitos.
		\item El codigo y el commit seria el siguiente:
	\end{itemize}	
	\begin{lstlisting}[language=bash,caption={Commit}][H]
		$ git commit -m "Creamos una clase Ejercito la cual nos va a poder ayudar al momento de querer datos de este como los soldados , el soldado de mayor vida promedio y tambien tenenemos en la clase Mapa la cual creamos objetos de esta clase Ejercito las cuales vamos a poder usar sus metodos"
	\end{lstlisting}	
	\begin{lstlisting}[language=java,caption={Las lineas de codigos de la clase Ejercito creada:}][H]
		import java.util.*;
		class Ejercito {   
			private ArrayList<ArrayList<Soldado>> army;
			private String kingdom;
			public Ejercito(){
			}
			public Ejercito(ArrayList<ArrayList<Soldado>> army, String kingdom){
				this.army = army;
				this.kingdom = kingdom;
			}
			public void viewSoldiers(String armyespe, int num, ArrayList<ArrayList<Soldado>> army){
				int numbersoldiers = 0;
				System.out.println("El Ejercito " + armyespe + " del " + num + " ejercito sus soldados son :");
				for(int i = 0; i < 10; i++){ //ITERACION
					for(int j = 0; j < 10 ; j++){//ITERACION
						if(army.get(i).get(j) != null){
							System.out.println("\n*********************************");
							System.out.println("El " + (numbersoldiers + 1) + " soldado es: ");
							System.out.println(army.get(i).get(j).toString());
							numbersoldiers++;
						}
					}
				}
			}
			public void longerLife(ArrayList<ArrayList<Soldado>> army, String kingdom){
				int mayor = 0;//METODO CREADO PARA PODER PERMITIRNOS A CONOCER EL SOLDADO CON MAYOR VIDA DE CADA EJERCITO 
				Soldado soldier = null;
				for(int i = 0; i < army.size(); i++){
					for(int j = 0; j < army.get(i).size(); j++){
						if(army.get(i).get(j) != null){ //COMPROBACION QUE HACEMOS PARA PODER DECIR QUE EL CASILLERO DONDE ESTAMOS ES UN SOLDADO QUE EXISTE
							if(army.get(i).get(j).getLifeActual() > mayor){ //COMPARAMOS PUNTOS DE VIDA DE CADA SOLDADO PARA VER QUIEN ES EL MAYOR 
								mayor = army.get(i).get(j).getLifeActual();
								soldier = army.get(i).get(j);
							}
						}
					}
				}
				System.out.println("El soldado con mayor vida del Ejercito " + kingdom + " es: ");
				System.out.println(soldier.toString());//IMPRIMIMOS SUS DATOS PARA PODER VER DE QUE SOLDADO SE TRATA 
				System.out.println("*********************************");
			}
			public double averageLife(ArrayList<ArrayList<Soldado>> army , String kingdom){
				int sum = 0;
				int count = 0;
				System.out.println("El promedio de puntos de vida del Ejercito " + kingdom + " es: ");
				for(int i = 0; i < 10; i++){
					for(int j = 0; j < 10; j++){ //ITERACION LA CUAL NOS AYUDA A PASAR POR TODOS LOS SOLDADOS DE CADA EJERCITO
						if(army.get(i).get(j) != null){ 
							sum += army.get(i).get(j).getLifeActual();
							count++;
						}
					}
				}
				if(sum != 0){
					double avg = sum / (count * 1.0);
					System.out.println(avg); // DAMOS A CONOCER EL PROMEDIO DE VIDA DE CADA EJERCITO 
					System.out.println("*********************************");
					return avg;
				}else{
					double avg = 0;
					System.out.println(avg); // DAMOS A CONOCER EL PROMEDIO DE VIDA DE CADA EJERCITO 
					System.out.println("*********************************");
					return avg;
				}
			}
			public void rankingInsercionLife(ArrayList<ArrayList<Soldado>> army , String kingdom){
				System.out.println("\nEl Ejercito " + kingdom + " ordenando por metodo insercion: ");
				int count = 0;
				for(int i = 0; i < 10; i++){
					for(int j = 0; j < 10; j++){ //ITERACION LA CUAL NOS AYUDA A PASAR POR TODOS LOS SOLDADOS DE CADA EJERCITO
						if(army.get(i).get(j) != null){ 
						   count++;
						}
					}
				}
				System.out.println("------------------------------------------");
				System.out.println("Mostrando Ranking del Ejercito " + kingdom + " ..... ////// --->");
				Soldado[] soldados = new Soldado[count];
				int x = 0;
				for(int i = 0; i < 10; i++){
					for(int j = 0; j < 10; j++){ //ITERACION LA CUAL NOS AYUDA A PASAR POR TODOS LOS SOLDADOS AL ARRAY SOLDADO PARA PODER USAR EL USO DEL METODO DE ORDENACION INSERCION
						if(army.get(i).get(j) != null){ 
							if(count - count + x == count){
								break;
							}else{
								soldados[count - count + x] = army.get(i).get(j); //LA MISMA LOGICA QUE EL ANTERIOR METODO SOLO QUE EN ESTE LO USARIAMOS DE MANERA DIFERENTE YA QUE ESTE SERIA DE FORMA DE INSERCION
							}
							x++;   
						}
					}
				}
				int n = soldados.length;
				for (int i = 1; i < n; i++) {
					Soldado actual = soldados[i];
					int j = i - 1;
					while (j >= 0 && soldados[j].getLifeActual() < actual.getLifeActual()) { //ORDENAMOS EL EJERCITO RESPECTIVAMENTE MEDIANTE EL METODO QUE NOS OFRECE INSERCION EL CUAL ES ESTE CODIGO
						soldados[j + 1] = soldados[j];
						j--;
					}
					soldados[j + 1] = actual;
				}
				for(int i = 0; i < soldados.length; i++){
					System.out.print("\n" + "Puesto " + (i + 1));
					System.out.println(soldados[i].toString()); //PUBLICAMOS RESULTADOS
					System.out.println("------------------");
				}
				System.out.println("*********************************");
			}
			public void setArmy(ArrayList<ArrayList<Soldado>> army){
				this.army = army;
			}
			public void setKingdom(String kingdom){
				this.kingdom = kingdom;
			}
			public String getKingdom(){
				return kingdom;
			}
		}
	\end{lstlisting}
	\begin{lstlisting}[language=java,caption={Las lineas de codigos de la clase Mapa creada:}][H]
		import java.text.DecimalFormat;
		import java.util.*;
		
		public class Mapa {
			Scanner sc = new Scanner(System.in);
			Random rdm = new Random();
			private String territory;
			private ArrayList<ArrayList<Soldado>> board;
			private Ejercito army1e;
			private Ejercito army2e;
			private String[] typesterritory = {"bosque", "campo abierto", "montana", "desierto", "playa"};
			private String[] kingdoms = {"Inglaterra", "Francia", "Sacro", "Castilla", "Aragon", "Moros"};
			public Mapa(){
				this.board = fillboard();
			}
			public void iniciarJuego() {
				menuBatalla();
				int n = sc.nextInt();
				while (n == 1) {
					String kingdom1 = kingdoms[rdm.nextInt(6)];
					String kingdom2 = kingdoms[rdm.nextInt(6)];
					ArrayList<ArrayList<Soldado>> army1 = fillarray(kingdom1, 1);
					ArrayList<ArrayList<Soldado>> army2 = fillarray(kingdom2, 2);
					army1e = new Ejercito(army1,kingdom1);
					army2e = new Ejercito(army2,kingdom2);
					territory = typesterritory[rdm.nextInt(5)];
					System.out.println("\n*********************************");
					System.out.println("El tipo de territorio es: " + territory);
					System.out.println("\n*********************************");
					bonificacion(army1, territory, kingdom1);
					bonificacion(army2, territory, kingdom2);
					army1e.viewSoldiers(kingdom1, 1, army1);
					army2e.viewSoldiers(kingdom2, 2, army2);
					viewBoard(army1, army2);
					army1e.longerLife(army1, kingdom1);
					army2e.longerLife(army2, kingdom2);
					double avg1 = army1e.averageLife(army1, kingdom1);
					double avg2 = army2e.averageLife(army2, kingdom2);
					army1e.rankingInsercionLife(army1, kingdom1);
					army2e.rankingInsercionLife(army2, kingdom2);
					viewBoard(army1, army2);
					resultBattleInfo(army1, kingdom1 , 1);
					resultBattleInfo(army2, kingdom2 , 2);
					int sum1 = resultBattleSum(army1, kingdom1, 1);
					int sum2 = resultBattleSum(army2, kingdom2, 2);
					double sumtotal = (sum1 * 1.0) + (sum2 * 1.0);
					resbattle(sumtotal, sum1, sum2, 1, 2, kingdom1, kingdom2);
					volveraJugar();
					n = sc.nextInt();
				}
			}
			public static ArrayList<ArrayList<Soldado>> fillboard(){
				ArrayList<ArrayList<Soldado>> army = new ArrayList<ArrayList<Soldado>>();
				for(int i = 0; i < 10; i++){ //ITERACION
					army.add(new ArrayList<Soldado>()); //LLENAMOS NUESTROS ARRAYLIST BIDIMENSIONAL CON CADA FILA PARA QUE CUMPLAN CON ESTRUCTURA DEL TABLERO
					for(int j = 0; j < 10 ; j++){//ITERACION
						army.get(i).add(null); // LLENAMOS CADA FILA DEL ARRAYLIST CON UN OBJETO SOLDADO CON TAL QUE ESTE SEA NULL PARA QUE SEPA QUE ESTE TIENE UNA CASILLA PERO NO HAY NADIE TODAVIA SE PUEDE LLENAR 
					}
				}
				return army;
			}
			public static ArrayList<ArrayList<Soldado>> fillarray(String armyespe, int num){
				Random rdm = new Random();
				ArrayList<ArrayList<Soldado>> army = new ArrayList<ArrayList<Soldado>>();
				int numbersoldiers = rdm.nextInt(10) + 1; //NUMERO DE SOLDADOS ALEATORIOS ENTRE 1 A 10 SOLDADOS 
				for(int i = 0; i < 10; i++){ //ITERACION
					army.add(new ArrayList<Soldado>()); //LLENAMOS NUESTROS ARRAYLIST BIDIMENSIONAL CON CADA FILA PARA QUE CUMPLAN CON ESTRUCTURA DEL TABLERO
					for(int j = 0; j < 10 ; j++){//ITERACION
						army.get(i).add(null); // LLENAMOS CADA FILA DEL ARRAYLIST CON UN OBJETO SOLDADO CON TAL QUE ESTE SEA NULL PARA QUE SEPA QUE ESTE TIENE UNA CASILLA PERO NO HAY NADIE TODAVIA SE PUEDE LLENAR 
					}
				}
				System.out.println("El Ejercito " + armyespe + " tiene " + numbersoldiers + " soldados : " ); 
				System.out.println("");
				for(int i = 0; i < numbersoldiers; i++){ //LLENAMOS CASILLAS CON CADA SOLDADO CREADO ALEATORIAMENTE
					Soldado soldado = getRandomSoldado();
					String name = "Soldado" + i + "X" + num;
					//System.out.println(name); PRUEBA QUE SE HIZO PARA VER LOS NOMBRES
					int health = rdm.nextInt(5) + 1;
					int row = rdm.nextInt(10) + 1;
					int speed = rdm.nextInt(5) + 1;
					String column = String.valueOf((char)(rdm.nextInt(10) + 65)); //REUTILIZAMOS CODIGO DEL ANTERIOR ARCHIVO VIDEOJUEGO2.JAVA YA QUE TENDRIAN LA MISMA FUNCIONALIDAD
					//System.out.println(army.get(row - 1).get((int)column.charAt(0) - 65)); PRUEBA QUE SE HIZO PARA COMPROBAR SI EL OBJETO SE ESTABA DANDO O NO CAPAZ NI EXISTIA  
					int lifelevel, defenselevel = 0;
					int attacklevel = 0;
					if (soldado instanceof Espadachin) {
						name = "Espadachin" + i + "X" + num; 
						lifelevel = rdm.nextInt(3) + 8; 
						attacklevel = 10;
						defenselevel = 8;
						soldado.setName(name);                  
						soldado.setAttackLevel(attacklevel);
						soldado.setDefenseLevel(defenselevel);                    
						soldado.setLifeLevel(lifelevel);
						soldado.setRow(row);
						soldado.setColumn(column);
						if(army.get(row - 1).get((int)column.charAt(0) - 65) == null){
							System.out.println("Registrando al " + (i + 1) + " soldado del Ejercito " + armyespe + "");
							army.get(row - 1).set((int)column.charAt(0) - 65, new Espadachin(name, attacklevel, defenselevel, lifelevel, speed, "Espadachin", true, row, column, attacklevel));
							army.get(row - 1).get((int)column.charAt(0) - 65).setSpeed(speed);
							System.out.println(army.get(row - 1).get((int)column.charAt(0) - 65).toString());
							System.out.println("---------------------------------");
						}else{
							i -= 1; //NOS AYUDARIA CON LOS SOLDADOS QUE SE REPITEN EN EL MISMO CASILLERO CON TAL QUE NO DEBERIA CONTAR 
						}
					} else if (soldado instanceof Arquero) {
						name = "Arquero" + i + "X" + num; 
						attacklevel = 7;
						defenselevel = 3;
						lifelevel = rdm.nextInt(3) + 3; 
						soldado.setName(name);                  
						soldado.setAttackLevel(attacklevel);
						soldado.setDefenseLevel(defenselevel);                    
						soldado.setLifeLevel(lifelevel);
						soldado.setRow(row);
						soldado.setColumn(column);
						if(army.get(row - 1).get((int)column.charAt(0) - 65) == null){
							System.out.println("Registrando al " + (i + 1) + " soldado del Ejercito " + armyespe + "");
							army.get(row - 1).set((int)column.charAt(0) - 65, new Arquero(name, attacklevel, defenselevel, lifelevel, speed, "Arquero", true, row, column, attacklevel));
							army.get(row - 1).get((int)column.charAt(0) - 65).setSpeed(speed);
							System.out.println(army.get(row - 1).get((int)column.charAt(0) - 65).toString());
							System.out.println("---------------------------------");
						}else{
							i -= 1; //NOS AYUDARIA CON LOS SOLDADOS QUE SE REPITEN EN EL MISMO CASILLERO CON TAL QUE NO DEBERIA CONTAR 
						}
					} else if (soldado instanceof Caballero) {
						name = "Caballero" + i + "X" + num; 
						attacklevel = 13;
						defenselevel = 7;
						lifelevel = rdm.nextInt(3) + 10; 
						soldado.setName(name);                  
						soldado.setAttackLevel(attacklevel);
						soldado.setDefenseLevel(defenselevel);                    
						soldado.setLifeLevel(lifelevel);
						soldado.setRow(row);
						soldado.setColumn(column);
						if(army.get(row - 1).get((int)column.charAt(0) - 65) == null){
							System.out.println("Registrando al " + (i + 1) + " soldado del Ejercito " + armyespe + "");
							army.get(row - 1).set((int)column.charAt(0) - 65, new Caballero(name, attacklevel, defenselevel, lifelevel, speed, "Caballero", true, row, column, false));
							army.get(row - 1).get((int)column.charAt(0) - 65).setSpeed(speed);
							System.out.println(army.get(row - 1).get((int)column.charAt(0) - 65).toString());
							System.out.println("---------------------------------");
						}else{
							i -= 1; //NOS AYUDARIA CON LOS SOLDADOS QUE SE REPITEN EN EL MISMO CASILLERO CON TAL QUE NO DEBERIA CONTAR 
						}
					} else if (soldado instanceof Lancero) {
						name = "Lancero" + i + "X" + num; 
						attacklevel = 5;
						defenselevel = 10;
						lifelevel = rdm.nextInt(3) + 5;
						soldado.setName(name);                  
						soldado.setAttackLevel(attacklevel);
						soldado.setDefenseLevel(defenselevel);                    
						soldado.setLifeLevel(lifelevel);
						soldado.setRow(row);
						soldado.setColumn(column);
						if(army.get(row - 1).get((int)column.charAt(0) - 65) == null){
							System.out.println("Registrando al " + (i + 1) + " soldado del Ejercito " + armyespe + "");
							army.get(row - 1).set((int)column.charAt(0) - 65, new Lancero(name, attacklevel, defenselevel, lifelevel, speed, "Lancero", true, row, column, attacklevel));
							army.get(row - 1).get((int)column.charAt(0) - 65).setSpeed(speed);
							System.out.println(army.get(row - 1).get((int)column.charAt(0) - 65).toString());
							System.out.println("---------------------------------");
						}else{
							i -= 1; //NOS AYUDARIA CON LOS SOLDADOS QUE SE REPITEN EN EL MISMO CASILLERO CON TAL QUE NO DEBERIA CONTAR 
						}
					}
				}
				System.out.println("*********************************");
				return army;
			}
			public static void menuBatalla(){
				System.out.println("-------------------------------------------");
				System.out.println("--                 MENU                  --"); 
				System.out.println("-------------------------------------------");
				System.out.println(" SELECCIONE UN NUMERO PARA PODER EMPEZAR O TERMINAR");
				System.out.println(" 1 : JUGAR");
				System.out.println(" 2 : NO JUGAR");
			}
			public static void viewBoard(ArrayList<ArrayList<Soldado>> army1, ArrayList<ArrayList<Soldado>> army2){ //EN ESTE METODO DEMOSTRAREMOS LA TABLA REUTILIZAREMOS CODIGOS DE ANTERIORES LABORATORIOS PARA PODER HACER LA BASE DE ESTE TABLERO
				System.out.println("\nMostrando tabla de posicion ... --");
				System.out.println("Leyenda: Ejercito1 --> 1#1 | Ejercito2 --> 2#2"); //RECONOCIMIENTO PARA LOS EJERCITOS Y POSICION DE SUS SOLDADOS
				System.out.println("\n \t   A\t   B\t   C\t   D\t   E\t   F\t   G\t   H\t   I\t   J"); // RECONOCIMIENTO PARA CADA UBICACION DE CADA SOLDADO EN EL TABLERO POR PARTE DE LAS COLUMNAS
				System.out.println("\t_________________________________________________________________________________");
				for(int i = 0; i < 10; i++ ){
					System.out.print((i + 1) + "\t"); // RECONOCIMIENTO PARA CADA UBICACION DE CADA SOLDADO EN EL TABLERO POR PARTE DE LAS FILAS
						for(int j = 0; j < 10; j++){
								if(army1.get(i).get(j) != null && army2.get(i).get(j) != null){ //CREAMOS UN IF PARA QUE ESTE NOS AYUDE A SABER QUIEN DE ESTOS SOLDADOS SE OCUPARA DEL CASILLERO EL CUAL DONDE ESTAN PELEANDO
									if(army1.get(i).get(j).getLifeActual() > army2.get(i).get(j).getLifeActual()){
										army1.get(i).get(j).setLifeActual(army1.get(i).get(j).getLifeActual() - army2.get(i).get(j).getLifeActual()); //Cambiamos 
										army2.get(i).set(j, null); 
										System.out.print("|  1" + obtenerInicial(army1.get(i).get(j)) + "1  ");
									}else if(army2.get(i).get(j).getLifeActual() > army1.get(i).get(j).getLifeActual()){
										army2.get(i).get(j).setLifeActual(army2.get(i).get(j).getLifeActual() - army1.get(i).get(j).getLifeActual());
										army1.get(i).set(j, null);;
										System.out.print("|  2" + obtenerInicial(army2.get(i).get(j)) + "2  ");
									}else{
										army2.get(i).set(j, null);
										army1.get(i).set(j, null);
										System.out.print("|   " + "" + "   ");
									}
								}else if(army1.get(i).get(j) != null){
									System.out.print("|  1" + obtenerInicial(army1.get(i).get(j)) + "1  ");
								}else if(army2.get(i).get(j) != null){
									System.out.print("|  2" + obtenerInicial(army2.get(i).get(j)) + "2  ");
								}else{
									System.out.print("|   " + " " + "   ");
								}
						}
						System.out.println("|");
						System.out.println("\t|_______|_______|_______|_______|_______|_______|_______|_______|_______|_______|");
				}
				System.out.println("\n*********************************");
			}
			public static Soldado getRandomSoldado() {
				Random rdm = new Random();
				int tipoSoldado = rdm.nextInt(4);
				switch (tipoSoldado) {
					case 0:
						return new Espadachin();
					case 1:
						return new Arquero();
					case 2:
						return new Lancero();
					case 3:
						return new Caballero();
					default:
						return new Espadachin();
				}
			}
			public static String obtenerInicial(Soldado soldado) {
				if (soldado instanceof Espadachin) {
					return "E";
				} else if (soldado instanceof Arquero) {
					return "A";
				} else if (soldado instanceof Caballero) {
					return "C";
				} else if (soldado instanceof Lancero) {
					return "L";
				} else {
					return "S";
				}
			}
			public void bonificacion(ArrayList<ArrayList<Soldado>> army, String territory , String kingdom) {
				for(int i = 0; i < 10; i++){ //ITERACION
					for(int j = 0; j < 10 ; j++){//ITERACION
						if(army.get(i).get(j) != null){
							if(kingdom.equals("Inglaterra") && territory.equals("bosque")){
								army.get(i).get(j).setLifeLevel(army.get(i).get(j).getLifeLevel() + 1);
							}else if(kingdom.equals("Francia") && territory.equals("campo abierto")){
								army.get(i).get(j).setLifeLevel(army.get(i).get(j).getLifeLevel() + 1);
							}else if((kingdom.equals("Castilla") || kingdom.equals("Aragon")) && territory.equals("montana")){
								army.get(i).get(j).setLifeLevel(army.get(i).get(j).getLifeLevel() + 1);
							}else if(kingdom.equals("Moros") && territory.equals("desierto")){
								army.get(i).get(j).setLifeLevel(army.get(i).get(j).getLifeLevel() + 1);
							}else if(kingdom.equals("Sacro") && (territory.equals("desierto") || territory.equals("playa") || territory.equals("campo abierto"))){
								army.get(i).get(j).setLifeLevel(army.get(i).get(j).getLifeLevel() + 1);
							}
						}
					}
				}
			} 
			public static void resultBattleInfo(ArrayList<ArrayList<Soldado>> army, String kingdom, int n){
				System.out.println("Ejercito " + n + " : " + kingdom);
				int numbersoldiers = 0;
				int numberespadachines = 0;
				int numbercaballeros = 0;
				int numberlanceros = 0;
				int numberarqueros = 0;
				for(int i = 0; i < 10; i++){ //ITERACION
					for(int j = 0; j < 10 ; j++){//ITERACION
						if(army.get(i).get(j) != null){
							numbersoldiers++;
							if(army.get(i).get(j) instanceof Espadachin){
								numberespadachines++;
							}else if(army.get(i).get(j) instanceof Caballero){
								numbercaballeros++;
							}else if(army.get(i).get(j) instanceof Lancero){
								numberlanceros++;
							}else if(army.get(i).get(j) instanceof Arquero){
								numberarqueros++;
							}
						}
					}
				}
				System.out.println("Cantidad total de soldados creados: " + numbersoldiers);
				System.out.println("Espadachines: " + numberespadachines);
				System.out.println("Arqueros: " + numberarqueros);
				System.out.println("Caballeros: " + numbercaballeros);
				System.out.println("Lanceros: " + numberlanceros + "\n");
			}
			public static int resultBattleSum(ArrayList<ArrayList<Soldado>> army, String kingdom, int n){
				int sum = 0;
				for(int i = 0; i < 10; i++){ //ITERACION
					for(int j = 0; j < 10 ; j++){//ITERACION
						if(army.get(i).get(j) != null){
							sum += army.get(i).get(j).getLifeLevel();
						}
					}
				}
				System.out.println("Ejercito " + n + ": " + kingdom + ": " + sum);
				return sum;
			}
			public static void resbattle(double sumtotal , int sum1, int sum2 , int n1, int n2 , String kingdom1, String kingdom2){
				DecimalFormat df = new DecimalFormat("#.##");
				String numero1 = df.format(((sum1 / sumtotal) * 100));
				String numero2 = df.format(((sum2 / sumtotal) * 100));
				if(sum1 > sum2){
					System.out.println("El ganador es el ejercito " + n1 + " de: " + kingdom1 +". Ya que al generar los porcentajes de probabilidad de victoria basada en los niveles de vida de sus soldados y aplicando un experimento aleatorio salio vencedor. (Aleatorio generado : "+ numero1 + ")");
				}else if (sum2 > sum1){
					System.out.println("El ganador es el ejercito " + n2 + " de: " + kingdom2 +". Ya que al generar los porcentajes de probabilidad de victoria basada en los niveles de vida de sus soldados y aplicando un experimento aleatorio salio vencedor. (Aleatorio generado : "+ numero2 + ")");
				}else{
					System.out.println("El resultado de la batalla es Empate");
				}
			}
			public static void volveraJugar(){
				System.out.println("\n*********************************");
				System.out.println(" DESEA VOLVER A JUGAR");
				System.out.println(" 1 : JUGAR");
				System.out.println(" 2 : NO JUGAR");
			}
		}
	\end{lstlisting}
	\subsection{Estructura de laboratorio 22}
	\begin{itemize}	
		\item El contenido que se entrega en este laboratorio22 es el siguiente:
	\end{itemize}
	\begin{lstlisting}[style=ascii-tree]
	\Lab22
	\end{lstlisting}    
	\section{\textcolor{red}{Rúbricas}}
	
	\subsection{\textcolor{red}{Entregable Informe}}
	\begin{table}[H]
		\caption{Tipo de Informe}
		\setlength{\tabcolsep}{0.5em} % for the horizontal padding
		{\renewcommand{\arraystretch}{1.5}% for the vertical padding
		\begin{tabular}{|p{3cm}|p{12cm}|}
			\hline
			\multicolumn{2}{|c|}{\textbf{\textcolor{red}{Informe}}}  \\
			\hline 
			\textbf{\textcolor{red}{Latex}} & \textcolor{blue}{El informe está en formato PDF desde Latex,  con un formato limpio (buena presentación) y facil de leer.}   \\ 
			\hline 
			
			
		\end{tabular}
	}
	\end{table}
	
	\clearpage
	
	\subsection{\textcolor{red}{Rúbrica para el contenido del Informe y demostración}}
	\begin{itemize}			
		\item El alumno debe marcar o dejar en blanco en celdas de la columna \textbf{Checklist} si cumplio con el ítem correspondiente.
		\item Si un alumno supera la fecha de entrega,  su calificación será sobre la nota mínima aprobada, siempre y cuando cumpla con todos lo items.
		\item El alumno debe autocalificarse en la columna \textbf{Estudiante} de acuerdo a la siguiente tabla:
	
		\begin{table}[ht]
			\caption{Niveles de desempeño}
			\begin{center}
			\begin{tabular}{ccccc}
    			\hline
    			 & \multicolumn{4}{c}{Nivel}\\
    			\cline{1-5}
    			\textbf{Puntos} & Insatisfactorio 25\%& En Proceso 50\% & Satisfactorio 75\% & Sobresaliente 100\%\\
    			\textbf{2.0}&0.5&1.0&1.5&2.0\\
    			\textbf{4.0}&1.0&2.0&3.0&4.0\\
    		\hline
			\end{tabular}
		\end{center}
	\end{table}	
	
	\end{itemize}
	
	\begin{table}[H]
		\caption{Rúbrica para contenido del Informe y demostración}
		\setlength{\tabcolsep}{0.5em} % for the horizontal padding
		{\renewcommand{\arraystretch}{1.5}% for the vertical padding
		%\begin{center}
		\begin{tabular}{|p{2.7cm}|p{7cm}|x{1.3cm}|p{1.2cm}|p{1.5cm}|p{1.1cm}|}
			\hline
    		\multicolumn{2}{|c|}{Contenido y demostración} & Puntos & Checklist & Estudiante & Profesor\\
			\hline
			\textbf{1. GitHub} & Hay enlace URL activo del directorio para el  laboratorio hacia su repositorio GitHub con código fuente terminado y fácil de revisar. &2 &X &2 & \\ 
			\hline
			\textbf{2. Commits} &  Hay capturas de pantalla de los commits más importantes con sus explicaciones detalladas. (El profesor puede preguntar para refrendar calificación). &4 &X &4 & \\ 
			\hline 
			\textbf{3. Código fuente} &  Hay porciones de código fuente importantes con numeración y explicaciones detalladas de sus funciones. &2 &X &2 & \\ 
			\hline 
			\textbf{4. Ejecución} & Se incluyen ejecuciones/pruebas del código fuente  explicadas gradualmente. &2 &X &2 & \\ 
			\hline			
			\textbf{5. Pregunta} & Se responde con completitud a la pregunta formulada en la tarea.  (El profesor puede preguntar para refrendar calificación).  &2 &X &2 & \\ 
			\hline	
			\textbf{6. Fechas} & Las fechas de modificación del código fuente estan dentro de los plazos de fecha de entrega establecidos. &2 &X &2 & \\ 
			\hline 
			\textbf{7. Ortografía} & El documento no muestra errores ortográficos. &2 &X &2 & \\ 
			\hline 
			\textbf{8. Madurez} & El Informe muestra de manera general una evolución de la madurez del código fuente,  explicaciones puntuales pero precisas y un acabado impecable.   (El profesor puede preguntar para refrendar calificación).  &4 &X &2 & \\ 
			\hline
			\multicolumn{2}{|c|}{\textbf{Total}} &20 & &18 & \\ 
			\hline
		\end{tabular}
		%\end{center}
		%\label{tab:multicol}
		}
	\end{table}
	
\clearpage

\section{Referencias}
\begin{itemize}			
	\item \url{https://drive.google.com/drive/folders/1_yO46UOaxs7uKVK7nrrkcNwybnk_ZJXJ}
\end{itemize}	
	
%\clearpage
%\bibliographystyle{apalike}
%\bibliographystyle{IEEEtranN}
%\bibliography{bibliography}
			
\end{document}